% II. Анализ на възможните решения и обосновка на избора.
% Разглеждат се четири решения: пространствен индекс, стратегия за векторизация,
% организация на конвейера и роля на външните библиотеки.
\section{Анализ на възможните решения и обосновка на избора}
\label{sec:alternatives}

\subsection{Пространствен индекс}

Три решения са приложими за търсене на съседи в облак от точки. Вокселната решетка
разделя пространството на кубчета с постоянен размер и намира съседите чрез обхождане
на съседните кубчета. Тя е тривиална за построяване и изключително бърза при равномерна
плътност, но при неравномерна плътност или силно издължена сцена или изразходва памет
за празни кубчета, или събира прекалено много точки в едно. K-d дървото~\cite{bentley1975kdtree}
се приспособява към плътността, защото разделя по медиана, и дава предсказуемо
поведение при заявка за $k$ най-близки съседи, но обхождането му е с нередовен достъп
до паметта. Октодървото стои между двете и се предпочита, когато е нужна работа на
няколко разделителни способности едновременно.

Изборът е k-d дърво като пространствен индекс, а вокселната решетка се използва в
другата си роля, като стъпка за прореждане преди индексирането. Първоначалният замисъл
беше решетката да бъде и второ, независимо изпълнение на същия интерфейс, за да се измери
едното срещу другото. От това е отстъпено: при неравномерна плътност, каквато има в
изпитваната сцена, решетката като индекс или изразходва памет за празни кубчета, или
събира прекалено много точки в едно, и измерването щеше да покаже това, което устройството
ѝ и без него казва. Приблизителните схеми~\cite{muja2014flann} са разгледани и оставени
извън обхвата, защото внасят втори източник на грешка, който би се смесил с грешката на
измервания алгоритъм.

\subsection{Стратегия за векторизация}

Ускоряването на заявките има три възможни пътя. Вътрешните функции на конкретния набор
инструкции дават пълен контрол, но обвързват изходния текст с една архитектура и налагат
отделна реализация за всяка. Библиотека за преносима векторизация запазва единен изходен
текст и оставя избора на набор инструкции на времето за компилация, но е външна
зависимост. Автоматичната векторизация на компилатора е най-евтина откъм усилие, но
резултатът зависи от компилатора и от неговата версия и не е част от договора на езика.

Избран е третият път, с една добавка, която премахва главния му недостатък. Вътрешният
цикъл е написан нарочно в пакетна форма: разстоянията за блок от точки се смятат в
отделен цикъл без разклонения, без обръщение към споделено състояние и с константен брой
итерации, а проверките са изнесени в следващ цикъл. Тази форма компилаторът разпознава.
Че я разпознава, не се предполага, а се проверява с \code{-fopt-info-vec}, и отговорът е
записан в Раздел~\ref{sec:implementation}.

Решаващият довод срещу библиотеката за преносима векторизация е ограничението за нулеви
задължителни зависимости, изложено по-долу. Условието изборът да работи е разположението
на данните: заявките се векторизират смислено само ако координатите са подредени по
компонента, а не по точка. Това превръща оформлението на паметта в проектно решение, а не
в подробност от реализацията, и е разгледано в Раздел~\ref{sec:design}.

\subsection{Организация на конвейера и роля на външните библиотеки}

Конвейерът може да бъде написан като монолитна функция, като верига от обекти с общ
интерфейс, или като съвкупност от независими преобразувания, свързвани при
конфигурирането. Първото е най-кратко, но прави отделната стъпка неизмерима.
Избрано е второто: всяка стъпка е обект с еднакъв договор, който приема облак и връща
облак заедно с показатели за изпълнението си. Така времето на стъпката е наблюдаемо
без промяна в кода, а редът на стъпките се задава отвън.

Разгледани бяха две положения за външните библиотеки. Първото е те да поемат входа и
изхода на данните и да служат за еталон при
сравнението~\cite{rusu2011pcl,bradski2000opencv}. Второто, избраното, е проектът да няма
нито една задължителна външна зависимост.

Доводът за второто е практически. Хранилището е публично, а изграждане, което изисква
мениджър на пакети, на практика не се изпълнява от никого при първия опит. Цената на
избора е обозрима и е платена: четците на PLY, PCD, PGM и PPM, рамката за проверка,
измервателната част и разлагането на матрица 3x3 са написани в проекта, общо под хиляда
реда. Загубеното е сравнението с чужда реализация. То е заместено с еталон, който не
изисква зависимост и е по-силен по своя въпрос: изчерпателното търсене, което е точно по
определение, и известното кораво преобразувание, срещу което се проверява съвместяването.

Алгоритмите, които са предмет на изследването, и при двете положения се реализират в
проекта. Обратното решение, тоест извикване на готовите реализации, би направило
измерването безсмислено, защото щеше да сравнява библиотека със самата себе си.
